% GENERATED FILE. Edit canonical lessons, then run npm run generate:books.
% canonical-source-hash: fnv1a64:bedee84a7041451d
% canonical-lessons: ZH-C416-yanwu, ZH-C416-chang, ZH-W416-chang-ground, ZH-C416-wuji, ZH-C416-zhuanji, ZH-C416-feixing

\chapter{A Storm}
\label{ch:zh-a-storm}
\hlchaptermodality{416}

\begin{chapteropening}
\textbf{By the end of this chapter:} \emph{I can say a flight is held up by a storm, that someone missed a plane, and talk about changing planes.}

Complete the last lesson of chapter 416: I can say a flight is held up by a storm, that someone missed a plane, and talk about changing planes.
\end{chapteropening}

\section[yánwù]{yánwù — to be delayed; to be held up}

\label{lesson:ZH-C416-yanwu}



\begin{quote}
Before the new one: say the Mandarin for the security check at an airport, then the Mandarin for a departure gate.
\end{quote}

\subsection*{You'll want to know: yánwù}
\emph{yánwù} — \textquotedblleft{}to be delayed; to be held up\textquotedblright{}.

\begin{quote}
\emph{hángbān yánwù le wǒmen yào zài jīchǎng zài děng liǎng ge xiǎoshí} — The flight's delayed: we have to wait at the airport two more hours.
\end{quote}

Its characters come later; for now it is a word for the ear.

\begin{cousinweb}
\emph{yánwù} is built from \emph{yán} 'to draw out' and \emph{wù} 'a mistake', both inherited from Middle Chinese, \emph{yen} and \emph{nguH} (Guangyun, in Baxter's transcription), which glosses the first 'to stretch, to lengthen' and the second 'an error'.
\end{cousinweb}

\begin{grammarlens}[title={What you've built}]
One more word for getting around.
\end{grammarlens}

\subsection*{Your turn}
\begin{itemize}
\raggedright
  \item \emph{Say it:} \emph{yánwù}
  \item \emph{Say it:} \emph{yánwù}, once more
  \item \emph{Say it:} \emph{hángbān yánwù le wǒmen yào zài jīchǎng zài děng liǎng ge xiǎoshí} — The flight's delayed: we have to wait at the airport two more hours.
  \item \emph{Recall:} say the Mandarin for the security check at an airport, then the Mandarin for a departure gate, then say \emph{yánwù} again
\end{itemize}

\begin{tcolorbox}[breakable,colback=teal!4,colframe=teal!35!black,title={Before you move on}]
What does yánwù mean? (\textquotedblleft{}To be delayed; to be held up\textquotedblright{}.) Say it once more.
\end{tcolorbox}
\section[\texorpdfstring{\zh{场} (chǎng)}{场 (chǎng)}]{\zh{场} (chǎng) — measure word for a rainstorm; measure word for a game or a show}

\label{lesson:ZH-C416-chang}



\begin{quote}
Before the new one: say the Mandarin for a departure gate, then the Mandarin for to be held up.
\end{quote}

\subsection*{You'll want to know: \zh{场}}
\textbf{\zh{场}} — \emph{chǎng} — \textquotedblleft{}measure word for a rainstorm; measure word for a game or a show\textquotedblright{}.

\begin{quote}
\textbf{\zh{今天有一场大雨}} \emph{jīntiān yǒu yì chǎng dà yǔ} — There's a big storm today.
\end{quote}

The new character, in pieces: \zh{场} — six strokes in two pieces: on the left the earth piece of \zh{地}, a bar, a vertical and a bar rising to the right; on the right a bar that turns down and zigzags to a hook, then two strokes falling to the left. The next lesson puts it in your hand.

\begin{cousinweb}
\textbf{\zh{场}} (\emph{chǎng}) is inherited from Middle Chinese \emph{drjang} (Guangyun, in Baxter's transcription), glossed as a place for offerings to the spirits and a threshing floor; how it came to count spells of rain and games is not settled. \emph{drjang} would regularly give a rising tone, cháng; the falling-rising tone of the measure word is irregular.
\end{cousinweb}

\begin{grammarlens}[title={What you've built}]
One more word for getting around.
\end{grammarlens}

\subsection*{Your turn}
\begin{itemize}
\raggedright
  \item \emph{Say it:} \emph{chǎng}
  \item \emph{Say it:} \emph{chǎng}, once more
  \item \emph{Say it:} \emph{jīntiān yǒu yì chǎng dà yǔ} — There's a big storm today.
  \item \emph{Recall:} say the Mandarin for a departure gate, then the Mandarin for to be held up, then say \emph{chǎng} again
\end{itemize}

\begin{tcolorbox}[breakable,colback=teal!4,colframe=teal!35!black,title={Before you move on}]
What does \zh{场} mean? (\textquotedblleft{}Measure word for a rainstorm; measure word for a game or a show\textquotedblright{}.) Say it once more.
\end{tcolorbox}
\section[\texorpdfstring{\zh{场} (chǎng)}{场 (chǎng)}]{\zh{场} — the earth piece, then a zigzag and two sweeps}

\label{lesson:ZH-W416-chang-ground}



\begin{quote}
The word first, then the character inside it.

\begin{itemize}
\raggedright
  \item \emph{Recall:} say \emph{chǎng} — \textbf{R1}, one lesson back, before it decays
\end{itemize}

\begin{quote}
\zh{地} \zh{机}
\end{quote}

\emph{Read it:} each one aloud, with nothing said first — \emph{dì}, \emph{jī}
\end{quote}

\subsection*{Script — the earth piece, then a zigzag and two sweeps}
\hlblockfigure{\detokenize{figures/ZH-W416-chang-ground-filmstrip.pdf}}{How \zh{场} is written, stroke by stroke}

\begin{quote}
\zh{场}
\end{quote}

\textbf{Two pieces, left before right.} Count the pieces, not the strokes.

\textbf{The left piece, three strokes}, as on the left of \zh{地}: a short bar, a vertical down through it, then a bar that rises to the right at the foot.

\textbf{The right piece, three strokes.} A bar to the right that turns down, zigzags and ends in a hook to the left: one stroke. Then two strokes falling to the left inside it, the second longer.

3 + 3 = six.

\subsection*{Your turn}
\begin{itemize}
\raggedright
  \item \emph{Write it:} \zh{场}, all six, with the model in front of you
  \item \emph{Cover:} write it again from memory, and count the pieces, not the strokes
\end{itemize}

\begin{tcolorbox}[breakable,colback=teal!4,colframe=teal!35!black,title={Before you move on}]
How many strokes in \zh{场}? (\textbf{Six}.) What is the left piece? (\textbf{The earth piece, as in \zh{地}}.) How many strokes falling to the left at the end? (\textbf{Two}.)

Source: \href{https://raw.githubusercontent.com/chanind/hanzi-writer-data/68d10a4b21150cae5e1ebbd223eed289cf32d90c/data/\%E5\%9C\%BA.json}{Hanzi Writer Data \zh{场}, pinned stroke paths and medians}.
\end{tcolorbox}
\section[wùjī]{wùjī — to miss a flight}

\label{lesson:ZH-C416-wuji}



\begin{quote}
Before the new one: say the Mandarin for to be held up, then the Mandarin for the measure word for a rainstorm.
\end{quote}

\subsection*{You'll want to know: wùjī}
\emph{wùjī} — \textquotedblleft{}to miss a flight\textquotedblright{}.

\begin{quote}
\emph{wǒmen zuò chūzūchē qù jīchǎng bùrán huì wùjī} — Let's take a taxi to the airport, or we'll miss the plane.
\end{quote}

Its characters come later; for now it is a word for the ear.

\begin{cousinweb}
\emph{wùjī} is built from \emph{wù} 'to miss' and \zh{机} 'a plane', both inherited from Middle Chinese, \emph{nguH} and \emph{kj+j} (Guangyun, in Baxter's transcription), which glosses the first 'an error'.
\end{cousinweb}

\begin{grammarlens}[title={What you've built}]
One more word for getting around.
\end{grammarlens}

\subsection*{Your turn}
\begin{itemize}
\raggedright
  \item \emph{Say it:} \emph{wùjī}
  \item \emph{Say it:} \emph{wùjī}, once more
  \item \emph{Say it:} \emph{wǒmen zuò chūzūchē qù jīchǎng bùrán huì wùjī} — Let's take a taxi to the airport, or we'll miss the plane.
  \item \emph{Recall:} say the Mandarin for to be held up, then the Mandarin for the measure word for a rainstorm, then say \emph{wùjī} again
\end{itemize}

\begin{tcolorbox}[breakable,colback=teal!4,colframe=teal!35!black,title={Before you move on}]
What does wùjī mean? (\textquotedblleft{}To miss a flight\textquotedblright{}.) Say it once more.
\end{tcolorbox}
\section[zhuǎnjī]{zhuǎnjī — to change planes}

\label{lesson:ZH-C416-zhuanji}



\begin{quote}
Before the new one: say the Mandarin for the measure word for a rainstorm, then the Mandarin for to miss a flight.
\end{quote}

\subsection*{You'll want to know: zhuǎnjī}
\emph{zhuǎnjī} — \textquotedblleft{}to change planes\textquotedblright{}.

\begin{quote}
\emph{wǒmen yào zài Běijīng zhuǎnjī} — We have to change planes in Beijing.
\end{quote}

Its characters come later; for now it is a word for the ear.

\begin{cousinweb}
\emph{zhuǎnjī} is built from \emph{zhuǎn} 'to turn' and \zh{机} 'a plane', both inherited from Middle Chinese, \emph{trjwenX} and \emph{kj+j} (Guangyun, in Baxter's transcription), which glosses the first 'to move, to turn'.
\end{cousinweb}

\begin{grammarlens}[title={What you've built}]
One more word for getting around.
\end{grammarlens}

\subsection*{Your turn}
\begin{itemize}
\raggedright
  \item \emph{Say it:} \emph{zhuǎnjī}
  \item \emph{Say it:} \emph{zhuǎnjī}, once more
  \item \emph{Say it:} \emph{wǒmen yào zài Běijīng zhuǎnjī} — We have to change planes in Beijing.
  \item \emph{Recall:} say the Mandarin for the measure word for a rainstorm, then the Mandarin for to miss a flight, then say \emph{zhuǎnjī} again
\end{itemize}

\begin{tcolorbox}[breakable,colback=teal!4,colframe=teal!35!black,title={Before you move on}]
What does zhuǎnjī mean? (\textquotedblleft{}To change planes\textquotedblright{}.) Say it once more.
\end{tcolorbox}
\section[fēixíng]{fēixíng — flight; flying (of an aircraft)}

\label{lesson:ZH-C416-feixing}



\begin{quote}
Before the new one: say the Mandarin for to miss a flight, then the Mandarin for to change planes on the way.
\end{quote}

\subsection*{You'll want to know: fēixíng}
\emph{fēixíng} — \textquotedblleft{}flight; flying (of an aircraft)\textquotedblright{}.

\begin{quote}
\emph{cóng zhèr zuò fēijī qù Běijīng fēixíng shíjiān shì liǎng ge xiǎoshí} — Flying from here to Beijing takes two hours.
\end{quote}

Its characters come later; for now it is a word for the ear.

\begin{cousinweb}
\emph{fēixíng} is built from \zh{飞} 'to fly' and \emph{xíng} 'to go', both inherited from Middle Chinese, \emph{pj+j} and \emph{haeng} (Guangyun, in Baxter's transcription), which glosses the second 'to walk, to go'.
\end{cousinweb}

\begin{grammarlens}[title={What you've built}]
One more word for getting around.
\end{grammarlens}

\subsection*{Your turn}
\begin{itemize}
\raggedright
  \item \emph{Say it:} \emph{fēixíng}
  \item \emph{Say it:} \emph{fēixíng}, once more
  \item \emph{Say it:} \emph{cóng zhèr zuò fēijī qù Běijīng fēixíng shíjiān shì liǎng ge xiǎoshí} — Flying from here to Beijing takes two hours.
  \item \emph{Recall:} say the Mandarin for to miss a flight, then the Mandarin for to change planes on the way, then say \emph{fēixíng} again
\end{itemize}

\begin{tcolorbox}[breakable,colback=teal!4,colframe=teal!35!black,title={Before you move on}]
What does fēixíng mean? (\textquotedblleft{}Flight; flying (of an aircraft)\textquotedblright{}.) Say it once more.
\end{tcolorbox}
